\def\bK{{\bb K}}

\bprob

    Let $G$ be a Lie group, and $H$ a closed embedded Lie subgroup.
    Define actions of $G$ on itself by left and right multiplication:
    $$ L,R\colon G\times G\to G,\qquad L(x,y)=xy,\quad R(x,y)=yx^{-1} . $$
    \benum
        \item Show that the action of $H$ on $G$ by $R$ is smooth, proper, and free.
        \item By the above point, we can form the quotient space $G/H$ with respect to the action $R$; let the quotient map be
        $\pi\colon G\to G/H$.
        Show that there is a unique $G$-action $\bar L$ on $G/H$ such that the following diagram commutes

        \commsq{G\times G&G\cr G\times G/H&G/H}
            {L}{\bar L}{\id\times\pi}{\pi}
    \eenum

\eprob

\benum
    \item First of all since $H$ is embedded, we can restrict $R$ to $H\times G\to G$, as shown in the previous homework.
    So the action is smooth.
    Clearly the action is free: if $R(x,y)=y$ then $yx^{-1}=y$ so $y=e$.
    To show that the action is proper, let $h_i\in H$ and $g_i\in G$ be sequences such that $g_i$ and $R(h_i,g_i)$ both converge; we need to show
    that $h_i$ has a convergent subsequence.
    Since $g_i$ and $R(h_i,g_i)=g_ih_i^{-1}$ converge, so does $R(h_i,g_i)^{-1}=h_ig_i^{-1}$ and thus $R(h_i,g_i)^{-1}g_i=h_i$ converges in $G$.
    Since $H$ is closed, it contains all of its limit points, thus $h_i$ converges in $H$, as desired.

    \item The orbit of $x\in G$ by $H$ under $R$ is
    $$ H\cdot x = {R(h,x)}[h\in H] = \set{xh^{-1}}[h\in H] = xH . $$
    If $\bar L$ exists it is unique, as it is given by
    $$ \bar L(g,xH) = \pi\circ L(g,x) = gxH . $$
    Note that this is well-defined: if $xH=yH$ then $gxH=gyH$.
    It is also a group action, since $\bar L(e,xH)=xH$ and $\bar L(g_1g_2,xH)=g_1g_2xH=\bar L(g_1,\bar L(g_2,xH))$.
    We need only show that it is smooth.

    Note that since $\id,\pi$ are surjective submersions, so is $\id\times\pi$ (the rank of the product of two maps is the sum of their ranks).
    Thus $\bar L$ is smooth if and only if $\bar L\circ(\id\times\pi)$ is.
    This is because if $M\xto\pi N$ is a smooth submersion, then $N\xto FP$ is smooth iff $F\circ\pi$ is smooth: one direction is clear,
    and the other follows by existence of local sections.
    If $F\circ\pi$ is smooth, let $p\in N$ and $q\in M$ such that $\pi(q)=p$.
    Then there is a local section $\sigma\colon U\to M$ with $q\in U$ and $\sigma(q)=p$, so
    $$ F|_U = F|_U\circ\id_U = F|_U\circ(\pi\circ\sigma) = (F|_U\circ\pi)\circ\sigma, $$
    which is smooth as the composition of smooth maps; so $F$ is smooth around every point in $N$, and thus smooth.

    Since $\bar L\circ(\id\times\pi)=\pi\circ L$ is smooth, $\bar L$ is smooth as desired.
\eenum

\bprob

    Let $G$ be a Lie group acting transitively on $M$.
    \benum
        \item Let $f\colon N\to Q$ be a smooth surjective map of constant rank.
        Show that $f$ is a submersion.
        \item For $p\in M$ show that $G_p\subseteq G$ is a closed embedded Lie subgroup.
        Conclude that $G/G_p$ is a smooth manifold.
        \item Let $\pi\colon G\to G/G_p$ be the quotient map, by the first question $G$ acts on $G/G_p$ by left multiplication.
        Show there is a unique $G$-equivariant diffeomorphism $F\colon G/G_p\to M$ such that $F(\pi(e))=p$.
        \item The unitary group $U(n)$ acts on $\bC^n$ by matrix multiplication, and thus it acts on the Grassmanian $G_\bC(k,n)$.
        Show that this action is smooth and transitive.
        Further show that the stabilizer of a point is the subgroup of block-diagonal matrices $U(k)\times U(n-k)$.
        Conclude that $G_\bC(k,n)$ is diffeomorphic $U(n)/(U(k)\times U(n-k))$.
    \eenum

\eprob

\benum
    \item Suppose $f$ has rank $k$, and let $n=\dim N,q=\dim Q$; suppose that $f$ is not a submersion so $k<q$.
    Then locally $f$ has the representation $\phi^{-1}\circ f\circ\psi\colon(x_1,\dots,x_n)\mapsto(x_1,\dots,x_k,0,\dots,0)$, where $\phi,\psi$ are coordinate charts
    of $Q,N$ respectively.
    This is a map between open subsets of Euclidean space, and since $f$ is surjective and $\phi,\psi$ are diffeomorphisms, it must be surjective.
    But its image consists of vectors whose final $q-k>0$ coordinates are zero, and this cannot be open in $\bR^q$ (any open ball centered on one of these points has points whose
    last coordinates are nonzero), a contradiction.
    Thus $f$ must have rank $k=q$, and is thus a submersion.

    \item Let $\theta\colon G\times M\to M$ be the group action, and define $\theta^p\colon G\to M$ by $\theta^p(g)=\theta(g,p)=gp$.
    Since $G$ acts on $M$ transitively, it is surjective.
    Furthermore it has constant rank since it is $G$-equivariant relative to $G$'s smooth action on itself by left multiplication:
    $$ \theta^p(g'g) = g'gp = g'\theta^p(g) . $$
    So by the previous point, it is a smooth submersion.

    Thus its level sets are embedded submanifolds, in particular $G_p=(\theta^p)^{-1}(p)$ is.
    It is also a subgroup of $G$, and since it is an embedded submanifold we can restrict the domains and codomains of multiplication and inversion to show that multiplication
    and inversion on it are smooth (as per the previous homework again).
    It is the preimage of a singleton under a continuous function, so it is closed as well, so is a closed embedded Lie subgroup.
    By the previous question, $G/G_p$ exists.

    \item If $F\colon G/G_p\to M$ is $G$-equivariant, it must satisfy
    $$ F(gG_p) = gF(G_p) = gF(\pi(e)) = gp $$
    for all $gG_p\in G/G_p$; so $F$ is unique if it exists.
    It is surjective since $G$ acts transitively on $M$, and injective since if $gp=g'p$ then $g^{-1}g'p=p$ so $g^{-1}g'\in G_p$ and thus $gG_p=g'G_p$.
    Now we need to show that it is smooth.

    As we showed in the previous question, $F$ is smooth iff $F\circ\pi$ is, and $F\circ\pi$ is the map $g\mapsto gp$, which is smooth as $G$ acts smoothly on $M$.
    Thus $F$ is a smooth $G$-equivariant bijection.
    Since it is $G$-equivariant, it has constant rank, and thus is a smooth submersion by our first point.
    Similarly it is an immersion: injective maps of constant rank are immersions by a similar argument to our first point (but even simpler, since $(x_1,\dots,x_n)\mapsto(x_1,\dots,x_k,0,\dots,0)$
    is injective only when $k=n$, we don't need to even consider topology).
    Also we could just note that $\dim G/G_p=\dim M$ so a smooth submersion is a smooth immersion.
    Then $F$ is a local diffeomorphism, again by the constant rank theorem, and since it is bijective it must be a global diffeomorphism.

    The first point that $F$ is a local diffeomorphism is because by the constant rank theorem it locally has the representation $\bar x\mapsto\bar x$, which is a diffeomorphism.
    Since it is a bijection, $F^{-1}$ exists, and for every point in its domain there is a neighborhood in which $F$ is a diffeomorphism, so $F^{-1}$ is smooth; thus $F^{-1}$
    is globally smooth.

    \item The action of $U(n)$ on $G_\bC(k,n)$ is given by $(U,V)\mapsto UV=\set{Uv}[v\in V]$.
    We will assume smoothness at first, and show it at the end.
    This action is transitive since if $V,V'\in G_\bC(k,n)$ are both subspaces of dimension $k$ then we can take orthonormal bases of them and define $U$ to map between them.
    Explicitly, let $B_0,B'_0$ be orthonormal bases of $V,V'$ respectively.
    Extend these to $B,B'$ orthonormal bases of $\bC^n$ and define $U$ to map $B$ to $B'$.
    This is a unitary matrix as it maps an orthonormal basis to an orthonormal basis, and since $UB_0=B'_0$, we have that $UV=V'$ as desired.

    For $V\in G_\bC(k,n)$, its stabilizer under $U(n)$ are all unitary matrices which map $V$ to itself.
    Such a unitary matrix must map $V^\perp$ to itself as well.
    Thus if $U$ is in the stabilizer, since $V\oplus V^\perp=\bC^n$, $U=U|_V\oplus U|_{V^\perp}$, so that $U$ is a block-diagonal matrix (up to equivalence).
    Since the restriction of $U$ to a $U$-invariant subspace remains unitary (it still maps orthonormal bases to orthonormal bases), $U|_V\in U(k)$ and $U|_{V^\perp}\in U(n-k)$
    so that $U\in U(k)\times U(n-k)$ (where $U(k)$ are the unitary transformations of $V$, and $U(n-k)$ the unitary transformations of $V^\perp$).
    Conversely any $(U_0,U_1)\in U(k)\times U(n-k)$ defines $U=U_0\oplus U_1$, a unitary matrix in $V$'s stabilizer.

    Thus by the previous point,
    $$ G_\bC(k,n) \cong U(n)/U(n)_V = U(n)/(U(k)\times U(n-k)) , $$
    as desired.

    Now to show that $U(n)$'s action on $G_\bC(k,n)$ is smooth, consider its charts $\Gamma_{V,W}\colon\hom(V,W)\to G_\bC(k,n)$ mapping $\ell$ to its graph.
    Recall that $\Gamma(\hom(V,W))$ is the collection of $k$-dimensional subspaces with trivial intersection with $W$.
    For $U\in U(n)$ and $\ell\in\hom(V,W)$, then
    $$ U\cdot\Gamma_{V,W}(\ell) = \set{Uv+U\ell v}[v\in V] , $$
    which has trivial intersection with $UW$, since $UV\oplus UW=\bC^n$, so it is in $\Gamma(\hom(UV,UW))$, and then
    $$ \Gamma^{-1}_{UV,UW}\circ U\circ\Gamma_{V,W}(\ell) = U\circ\ell\circ U^{-1}, $$
    since its graph is $\set{Uv+U\ell U^{-1}Uv}[v\in V]=\set{Uv+U\ell v}[v\in V]=U\Gamma(\ell)$.
    This is a smooth map, since conjugation is smooth ($\gl_n(\bC)$, and thus $U(n)$ since it is embedded, acts smoothly on $\mat_n(\bC)$).

    So each coordinate representation of the action is smooth, but we need to show that for each $(U,V)$, they lie in a neighborhood which is mapped into $\Gamma(\hom(UV,UW))$.
    Consider the map $U(n)\times\hom(V,W)\to\hom(V,\bC^n)$ which maps $(A,\ell)\mapsto A\ell$.
    This is smooth, as it is just matrix multiplication.
    Since $V$ is just represented by $0\in\hom(V,W)$ (meaning $\Gamma_{V,W}(0)=V$), we want to show that the preimage of the maps whose graphs don't intersect $UW$ is open.
    Now $\ell\in\hom(V,\bC^n)$'s graph doesn't intersect $UW$ if for a basis $v_1,\dots,v_k\in V$ and a basis $w_1,\dots,w_{n-k}\in W$, the collection $v_i+\ell v_i$ with $Uw_j$ remains
    linearly independent.
    That is, the determinant of $(\dots,v_i+\ell v_i,\dots,Uw_j,\dots)$ is nonzero.
    This is an open set, since mapping $\ell$ to this matrix is linear and thus continuous, and the determinant is continuous.

    If we let $O\subseteq U(n)\times\hom(V,W)$ be this open set, so consisting of pairs $(A,\ell)$ such that $\Gamma(A\ell)\cap W=0$, then through $\id\times\Gamma_{V,W}$ this gives an open
    set $O$ in $G_\bC(k,n)$.
    And since $0\in\hom(V,W)$ represents $V$, and $UV\cap UW=0$, this is an open neighborhood of $(U,V)$.
    Thus for every $(U,V)\in U(n)\times G_\bC(k,n)$ there is an open neighborhood $O$ which is mapped into an open chart $UV\in\Gamma(\hom(UV,UW))$, and as we showed the coordinate
    representation is smooth.
    Thus $U(n)$'s action on $G_\bC(k,n)$ is smooth
\eenum

\bprob

    Let $M$ be a smooth manifold, and $H$ a Lie group which acts on it smoothly, freely, and properly by $\phi\colon H\times M\to M$.
    Let $g$ be an $H$-invariant metric on $M$, and denote by $\pi\colon M\to M/H$ the projection.
    \benum
        \item Prove there exists a unique metric $\bar g$ on $M/H$ such that $\pi$ is a Riemannian submersion.
        \item Let $M=\bR^n$ and $H=\bZ^n$, and $\phi$ given by translation, so $M/H=T^n$.
        Let $g=g_0$ be the Euclidean metric.
        Construct an isometric embedding $(T^n,\bar g)\to(\bR^{2n},g_0)$.
    \eenum

\eprob

\benum
    \item For $p\in M$, we can write $d\phi_p$ by
    $$ d\pi_p\colon T_p(Hp)\oplus T_p(Hp)^\perp = T_pM \to T_{Hp}(M/H) , $$
    and since $T_pHp=\ker d\pi_p$ and $\pi$ is a submersion, we have that $d\pi_p$ forms an isomorphism $T_p(Hp)^\perp\to T_{Hp}(M/H)$, let us denote this by $d\hat\pi_p$.
    In order for $\pi$ to be a Riemannian submersion, we must then define
    $$ \bar g(v,u) = g(d\hat\pi_p^{-1}v,d\hat\pi_p^{-1}u),\qquad v,u\in T_{Hp}(M/H), $$
    for $v,u\in T_{Hp}(M/H)$.
    Thus $\bar g$ is unique if it exists.
    Since $g$ is a Riemannian metric, this is a metric at each point (in general $\iprod{T\bul,T\bul}$ is an inner product when $T$ is invertible).

    Note that this is not necessarily well-defined: if $Hp=Hq$ then this definition may give different values if $d\hat\pi_p$ or $d\hat\pi_q$ is used.
    But since $g$ is $H$-invariant, if $Hp=Hq$ suppose $q=\phi_h(q)$ then since $\pi\circ\phi_h=\pi$ we have that $d\pi_{p}\circ d(\phi_h)_q=d\pi_q$ so by $H$-invariance,
    $$ g(d\hat\pi_p^{-1}v,d\hat\pi_p^{-1}u) = g(d(\phi_{h^{-1}})d(\hat\pi_p)^{-1}v,d(\phi_{h^{-1}})d(\hat\pi_p)^{-1}u) = g(d\hat\pi_q^{-1}v,d\hat\pi_q^{-1}u) . $$

    Now we want to show that $\bar g$ is smooth.
    Notice that we can apply Gram-Schmidt to a smooth local coordinate frame to obtain a smooth orthonormal local coordinate frame (smoothness is due to $g$'s smoothness).
    Then for $p\in M$ we can define $\Pi_p\colon T_pM\to T_pM$ which is the orthonormal projection onto $T_p(Hp)^\perp$.
    Considering a smooth orthonormal coordinate frame, we see that this maps smooth vector fields to smooth vector fields, and furthermore it is clearly $C^\infty$-linear,
    and thus $\Pi\colon\X(M)\to\X(M)$ forms a $(1,1)$-tensor.
    Since $\pi$ is a submersion, it has local sections: there is a neighborhood $p\in U$ and a smooth local section $\sigma\colon\pi(U)\to M$ with $\sigma(Hp)=p$.
    Now we claim that locally, $d\hat\pi^{-1}=\Pi\circ d\sigma$.
    Since for $q\in U$, $\Pi\circ d\sigma$ maps into $T_q(Hq)^\perp$, we need only show that $d\hat\pi\circ\Pi\circ d\sigma=\id$.
    Since $d\pi_q$ is linear and $T_q(Hq)$ is its kernel as well as the projection $\Pi_q$'s, we see that $d\hat\pi_q\circ\Pi=d\pi_q$, and so
    $$ d\hat\pi_q\circ\Pi_q\circ d\sigma_{Hq} = d\pi_q\circ d\sigma_{Hq} = \id . $$
    Thus $d\hat\pi^{-1}=\Pi\circ d\sigma$ on $\pi(U)$, and thus locally
    $$ \bar g(v,u) = g(\Pi\circ d\sigma(v),\Pi\circ d\sigma(u)) , $$
    and since both $\Pi,d\sigma$ are smooth, so is $\bar g$ (locally around every point, and thus globally).

    \item For $n=1$, the obvious embedding $\bR/\bZ=T^1\to\bR^2$ is given by $\theta\mapsto(\cos(2\pi\theta),\sin(2\pi\theta))$, or under the natural identification $\bR^2\cong\bC$,
    $\theta\mapsto\exp(2\pi i\theta)$.
    This is smooth since precomposing it with the projection $\bR\to\bR/\bZ$ maps $x\in\bR$ to $\exp(2\pi ix)$ (since $\exp(2\pi ix)$ has a periodicity of $1$), which is smooth.
    This extends to an obvious embedding $T^n\to\bR^{2n}\cong\bC^n$ by
    $$ F\colon (\theta_1,\dots,\theta_n) \mapsto (\exp(2\pi i\theta_1),\dots,\exp(2\pi i\theta_n)) , $$
    which is again smooth.

    Note that $\pi\colon\bR^n\to T^n$ is a local diffeomorphism, as it is a submersion and $\dim T^n=\dim\bR^n=n$.
    Thus let $\ipdv{\theta^i}$ be the pushforward of the coordinate vector fields $\ipdv{x^i}$, and since $x^i$ are orthonormal and $\pi$ is a local Riemannian diffeomorphism (by construction it is a
    Riemannian submersion), $\theta^i$ are orthonormal as well.
    Then for $\bar\theta=(\theta_1,\dots,\theta_n)\in T^n$, (using the Jacobian of $F$ by $x^i$),
    $$ dF_{\bar\theta}\pdv{\theta^i} = -2\pi\sin(2\pi\theta_i)\pdv{x^i_1} + 2\pi\cos(2\pi\theta_i)\pdv{x^i_2} , $$
    where the global coordinates of $\bR^{2n}$ are labeled $x^1_1,x^1_2,\dots,x^n_1,x^n_2$.
    Then $dF_{\bar\theta}\ipdv{\theta^i}$ are orthogonal (since $\pdv{x^i_1},\pdv{x^i_2}$ are), and
    $$ g_0\parens{dF_{\bar\theta}\pdv{\theta^i},dF_{\bar\theta}\pdv{\theta^i}} = 4\pi^2\sin(2\pi\theta_i)^2 + 4\pi^2\cos(2\pi\theta_i)^2 = 4\pi^2 . $$
    So if we define $F'=F/2\pi$, the above computations normalize to give $1$.
    So $dF'$ maps an orthonormal basis to an orthonormal set, and is thus a Riemannian embedding.
\eenum

\bprob

    On $M=\bR^2$ consider the map
    $$ \phi\colon(0,\infty)\times(0,2\pi)\to M,\qquad \phi(r,\theta) = (r\cos\theta,r\sin\theta) . $$
    \benum
        \item Show that $\phi$ is a coordinate chart for $M$.
        \item Let $\set{\partial_r,\partial_\theta}$ denote the coordinate vector fields.
        For $a,b,c,d\in\bR$ define the vector field $X$ on $\phi$'s coordinate domain:
        $$ X = r(a\cos^2\theta+(b+c)\sin\theta\cos\theta+d\sin^2\theta)\partial_r + (c\cos^2\theta+(d-a)\sin\theta\cos\theta-b\sin^2\theta)\partial_\theta . $$
        Show that $X$ extends smoothly to a vector field on all of $M$.
        \item For any $p\in M$, calculate the integral curve of $X$ starting at $p$.
        Show that $X$ admits global flow.
        \item Draw the integral curves of $X$ for the following values of $a,b,c,d$.

        \medskip
        \centerline{\vbox{\halign{\strut\hfil#\hfil\kern.25cm\vrule\tabskip=.5cm&\hfil$#$&\hfil$#$&\hfil$#$&\hfil$#$\cr
            Case&a&b&c&d\cr\noalign{\hrule}
            Sink&-1&0&0&-2\cr
            Source&1&0&0&2\cr
            Saddle&1&0&0&-1\cr
            Spiral sink&-1&-1&1&-1\cr
            Spiral source&1&-1&1&1\cr
            Center&0&-1&1&0\cr
            Improper sink&-1&1&0&-1\cr
            Improper source&1&1&0&1\cr
        }}}
    \eenum

\eprob

\def\atan{{\rm atan}}

\benum
    \item $\phi$'s image is the set of all points $(x,y)$ where either $y\neq0$ or $x<0$ ($\bR^2$ without $[0,\infty)\times\bR$), which is open.
    It is smooth, as its components are smooth, and it is injective.
    Furthermore, its Jacobian is
    $$ J\phi = \pmatrix{\cos\theta&-r\sin\theta\cr \sin\theta&r\cos\theta} , $$
    which has determinant $r\neq0$, so it is invertible everywhere.
    Thus $\phi$ is a diffeomorphism onto its image, meaning it is a chart.

    \item Let us write $X=X^r\partial_r+X^\theta\partial_\theta$.
    We know that
    $$ \eqalign{
        \partial_r &= \pdvof xr\pdv x + \pdvof yr\pdv y = \cos\theta\pdv x + \sin\theta\pdv y\cr
        \partial_\theta &= \pdv x\theta\pdv x + \pdvof y\theta\pdv y = -r\sin\theta\pdv x + r\cos\theta\pdv y .
    } $$
    Thus we see that $X=X^1\ipdv x+X^2\ipdv y$, where
    $$ X^1 = X^r\cos\theta - rX^\theta\sin\theta,\qquad X^2 = X^r\sin\theta + rX^\theta\cos\theta . $$
    We have that
    $$ x = r\cos\theta,\quad y = r\sin\theta,\quad r^2 = x^2+y^2,\quad \cos^2\theta = \frac{x^2}{r^2},\quad \cos\theta\sin\theta = \frac{xy}{r^2},\quad \sin^2\theta = \frac{y^2}{r^2} . $$
    Thus some tedious calculations yield
    $$ X = (ax+by)\pdv x + (cx+dy)\pdv y . $$
    This can be clearly extended smoothly to all of $M$.

    \item Since $X$ has the form
    $$ X(x,y) = \pmatrix{a&b\cr c&d}\pmatrix{x\cr y} = A \pmatrix{x\cr y}, $$
    solving $\gamma'(t)=X(\gamma(t))$ boils down to the linear system of ODEs:
    $$ \pmatrix{\gamma_1'(t)\cr \gamma_2'(t)} = A\pmatrix{\gamma_1(t)\cr\gamma_2(t)} . $$

    Since $A$ is two-dimensional, it has a Jordan normal form which is diagonal, or is a $2\times2$ Jordan block.
    Suppose that $P^{-1}AP=J$ is $A$'s Jordan form, then a solution to $\alpha'=J\alpha$ gives a solution $\gamma=P\alpha$ to $\gamma'=A\gamma$, since
    $$ \gamma' = P\alpha' = PJ\alpha = AP\alpha = A\gamma , $$
    and vice versa.
    Note that $\alpha(0)=P^{-1}\gamma(0)=P^{-1}p$.

    If $A$ is diagonalizable then say $J=\diag(\lambda_1,\lambda_2)$.
    Solving for $\alpha$ is easy, as the ODE is $\alpha_i'=\lambda_i\alpha_i$ for $i=1,2$, and solving this gives $\alpha_i=c_i\exp(\lambda_it)$, where $c=P^{-1}p$.

    Now if $A$'s Jordan form is a Jordan block $J=J_2(\lambda)$, then the ODE becomes mildly more difficult:
    $$ \alpha_1' = \lambda\alpha_1 + \alpha'_2,\qquad \alpha_2' = \lambda\alpha_2 . $$
    This gives us $\alpha_2=c_2\exp(\lambda t)$ as before, and then $\alpha_1=c_1\exp(\lambda t)+c_2t\exp(\lambda t)$ (using the usual method for solving a non-homogeneous linear ODE).

    Note that $\lambda_1,\lambda_2,\lambda$ may be complex numbers.
    In this case, the real and imaginary parts of the solution give valid solutions to the real linear system of ODEs.
    For example, if $\lambda_1=\lambda_2=1+i$ then our solutions become
    $$ \gamma_{1,2} = \exp((1+i)t) = p_{1,2}e^t(\cos t+i\sin t) , $$
    and this gives us two solutions:
    $$ \gamma_{1,2} = e^t\cos t,\quad\hbox{and}\quad \gamma_{1,2} = e^t\sin t . $$
    Then the general solution is of the form
    $$ \gamma_{1,2} = Ae^t\cos t+Be^t\sin t . $$

    In either case $\alpha$, and thus $\gamma$, are defined on all of $\bR$ and thus $X$ has global flow.

    \item
    \benum
        \item Sink: the matrix is already diagonal and the solution becomes
        $$ \gamma(t) = \pmatrix{A\exp(-t)\cr B\exp(-2t)} . $$
        Thus the flow in general is given by curves (parabolas) of the form $(s,as^2)$ (pointing toward the origin, as when $t\to\infty$, $\exp(-t)\to0$).

        \item Source: the solution becomes $\gamma(t)=(A\exp(t),B\exp(t))$, and so the flow is given by curves of the form $(s,as^2)$ pointing away from the origin.

        \item Saddle: the solution becomes $\gamma(t)=(A\exp(t),B\exp(-t))$, so the flow is given by curves of the form $(s,as^{-1})$ pointing away from the origin.

        \item Spiral sink: the matrix is diagonalizable to $\diag(-1+i,-1-i)$ with eigenvectors $(1,-i)$ and $(1,i)$.
        The solution to the diagonalized problem is $\alpha_1=A\exp((-1+i)t)$ and $\alpha_2=B\exp((-1-i)t)$.
        But notice that since our ODE is a homogeneous linear system of order $2$, we need only find $2$ linearly independent solutions.
        Furthermore, given a complex solution its real and imaginary parts are both solutions to the real problem.
        Thus we can consider only $\alpha=(\exp(-1+i)t),0)$ (i.e.\ $A=1,B=0$), which gives us the complex solution
        $$ \gamma(t) = P\alpha = \pmatrix{1&1\cr -i&i}\pmatrix{\exp((-1+i)t)\cr0} = \pmatrix{\exp((-1+i)t)\cr -i\exp((-1+i)t)} . $$
        The real and complex parts of this are
        $$ \gamma_1(t) = e^{-t}\pmatrix{\cos t\cr \sin t},\qquad \gamma_2(t) = e^{-t}\pmatrix{\sin t\cr -\cos t} . $$
        So the general flow lines are given by
        $$ \gamma(t) = A\gamma_1(t) + B\gamma_2(t) . $$
        This forms a swirl with arrows pointing toward the origin (since as $t\to\infty$, the curve goes to $0$).

        \item Spiral source: the matrix is diagonalizable to $\diag(1+i,1-i)$ with eigenvectors $(1,-i)$ and $(1,i)$.
        As above, we have a complex solution
        $$ \pmatrix{\exp((1+i)t)\cr -i\exp((1+i)t)} , $$
        whose real and complex parts are
        $$ \gamma_1(t) = e^t\pmatrix{\cos t\cr \sin t},\qquad \gamma_2(t) = e^t\pmatrix{\sin t\cr -\cos t} . $$
        Thus the flow lines are just linear combinations of these, which look similar to the spiral sink but with arrows pointing out away from the origin.

        \item Center: the matrix is diagonalizable to $\diag(i,-i)$ with eigenvectors $(1,-i)$ and $(1,i)$.
        We have a complex solution
        $$ \pmatrix{\exp(it)\cr -i\exp(it)} , $$
        which has real and complex parts
        $$ \gamma_1(t) = \pmatrix{\cos t\cr \sin t},\qquad \gamma_2(t) = \pmatrix{\sin t\cr -\cos t} . $$
        The general flow lines are of the form
        $$ \gamma(t) = \pmatrix{A\cos t + B\sin t\cr A\sin t - B\cos t} , $$
        which are circles (with radius squared $A^2+B^2$, going counterclockwise.

        \item Improper sink: this matrix is in Jordan form, and so the solution is (as proven before):
        $$ \gamma(t) = e^{-t}\pmatrix{A + tB\cr B} . $$
        These form lopsided spirals into the origin.

        \item Improper source: this matrix is in Jordan form, and so the solution is:
        $$ \gamma(t) = e^t\pmatrix{A + tB\cr B} . $$
        This is similar to the improper sink, but the spirals go to infinity (and are mirrored along the $y$-axis).
    \eenum
\eenum



